\documentclass[10pt,a4paper]{amsart}
\usepackage[margin=19mm]{geometry}
\usepackage{amsmath,amssymb,mathtools}
\usepackage[scr=rsfs,cal=euler]{mathalfa}
\usepackage{xfrac}
\usepackage[unicode,hidelinks,bookmarks=false]{hyperref}
\newcommand{\ie}{i.e.\ }
\newcommand{\cf}{cf.\ }
\newcommand{\resp}{respectively\ }
\newcommand{\Subsection}{\subsection}
\providecommand{\nobreakdash}{\nobreak\hbox{-}}
\setlength{\emergencystretch}{2em}
\multlinegap0pt
\usepackage{amssymb}
\usepackage{enumerate}
\usepackage{xcolor}
\newtheorem{proposition}{Proposition}
\newcommand{\CI}{{\mathcal I}}
\DeclareMathOperator*{\E}{\mathbb{E}}
\newcommand{\N}{{\mathbb N}}
\newcommand{\abs}[1]{\mathopen{}\left| #1\mathclose{}\right|}
\newcommand{\be}{{\mathbf{e}}}
\newcommand{\brac}[1]{\mathopen{}\left( #1 \mathclose{}\right)}
\newcommand{\nnorm}[1]{\lvert\!|\!| #1|\!|\!\rvert}
\newcommand{\norm}[1]{\left\Vert #1\right\Vert}
\title{Joint ergodicity - 40 years on}
\author{Borys Kuca}
\date{Source: arXiv:2603.18974v1 (CC BY 4.0)}
\begin{document}
\maketitle

\noindent\emph{Source and licence.} Joint ergodicity - 40 years on. Version arXiv:2603.18974v1 (CC BY 4.0), distributed by arXiv under the Creative Commons Attribution 4.0 International licence, http://creativecommons.org/licenses/by/4.0/, as recorded in the arXiv licence metadata for this exact version and shown on the abstract page https://arxiv.org/abs/2603.18974v1. Full licence notice: \texttt{licenses/arxiv-2603.18974-CC-BY-4.0.txt}.\par\smallskip

\noindent\textit{Editorial scope.} Counted unit: the article's proposition s=3 (source0.tex lines 2676--2679). The article's complete original proof of it is delivered with it (source0.tex lines 2681--2811, one \texttt{proof} environment of 131 lines). No mathematics is written, repaired or re-derived: the statement and the proof are the article's own lines, in the article's own notation, and no retained line is substituted or re-flowed.  Retained uncounted as the source-local context the proof needs: lines 2579--2582; lines 2602--2604.  Nothing was omitted from any retained span, so the package carries no omission record.  No retained line contains a citation, so the wrapper carries no bibliography and no bibliography entry.  No retained line contains a figure environment, an image-inclusion command or drawing code, so no image is delivered and the package declares no asset dependency.  Editorial formatting only, in addition: an AMS \texttt{amsart} preamble that restates the article's own declarations and macros used by the retained lines, namely the environments \texttt{proposition} on separate counters, together with the standard packages those lines need.  The retained environments print in this document as Proposition 1 in order of appearance, whereas the article prints the same items with the numbering of its own full document.  The complete native source of the article as distributed in the arXiv e-print for this exact version is bundled unchanged as \texttt{sources/196743/source0.tex}.
\begin{proposition}[Degree lowering, baby case I]\label{P: degree lowering s=2}
    Assume %that $T_1, T_2$ are ergodic, and 
    that Property $(P_1)$ holds. If $\nnorm{F_2}_{2, T_2}>0$, then also $\nnorm{F_2}_{1, T_2}>0$. 
\end{proposition}

    \begin{align}\label{E: expanding dual}
        \lim_{N\to\infty}\E_{n\in[N]}\int T_2^{-a_2(n)}\overline{f_0}\cdot T_1^{a_1(n)}T_2^{-a_2(n)}\overline{f_1}\cdot \overline{\chi_2}\cdot 1_U\; d\mu > 0.
    \end{align}

\begin{proposition}[Degree lowering, baby case II]\label{P: degree lowering s=3}
    Assume %that $T_1, T_2$ are ergodic, and 
    that Property $(P_1)$ holds. If $\nnorm{F_2}_{3, T_2}>0$, then also $\nnorm{F_2}_{2, T_2}>0$. 
\end{proposition}

\begin{proof}
    First, we observe that the inverse theorem for degree-3 Host-Kra seminorm involves much more complicated objects than nonergodic eigenfunctions. To land in the situation where Property $(P_1)$ is applicable, we apply the inductive formula for Host-Kra seminorms. It gives
    %concerns eigenfunctions, we need to bring our assumption to a situation in which 
    \begin{align*}
        \lim_{H\to\infty}\E_{h\in[H]}\nnorm{\Delta_{h\be_2} F_2}_{2, T_2}^4 >0,
    \end{align*}
    where $\Delta_{h\be_2}f:=f\cdot T_2^h\overline{f}$.
    %  By the inverse theorem for the degree-2 Host-Kra seminorm, for every $h\in\N$ there exists an eigenfunction $\chi_{2,h}$ of $T_2$ such that 
    % \begin{align*}
    %     \lim_{H\to\infty}\E_{h\in[H]}\int \Delta_{h\be_2} F_2 \cdot \overline{\chi_{2,h}}\; d\mu >0,
    % \end{align*}     
    % and furthermore
    % \begin{align}\label{E: condition on positivity}       
    % \end{align}
    In particular, there exist $\delta>0$ and  a set $\Lambda\subseteq \N$ of positive lower density such that 
    \begin{align}\label{E: lower bound on Lambda}
        \nnorm{\Delta_{h\be_2} F_2}_{2, T_2}^4 \geq \delta \quad \textrm{for all}\quad h\in \Lambda.
    \end{align}
    By the inverse theorem for the degree-2 Host-Kra seminorm, for every $h\in\Lambda$ there exists an eigenfunction $\chi_{2,h}$ of $T_2$ such that 
    \begin{align*}%\label{E: positivity s=2}
        \int \Delta_{h\be_2} F_2 \cdot {\chi_{2,h}}\; d\mu \geq \delta\quad \textrm{and} \quad \E(\Delta_{h\be_2} F_2\cdot {\chi_{2,h}}|\CI(T_2))\geq 0.
    \end{align*}
    We will show once again that the nonergodic eigenfunctions $\chi_{2,h}$ admit some special structure, and specifically that on positive-measure sets, they take the form $\chi\cdot \lambda_h$ for some nonergodic eigenfunctions $\chi$ of $T_2$ and $T_2$-invariant $\lambda_h$'s. (If $T_2$ is ergodic, then simply $\chi_{2,h} = \chi$ for $h\in\Lambda$.)
    %(for $h\in\Lambda$)
    
    By the popularity principle, for every $h\in\Lambda$ we can find a $T_2$-invariant set $U_h$ with $\mu(U_h)\geq \delta/2$ such that
    %Consequently, we can find $\delta > 0$ and a $T$-invariant set $U\subseteq X$ so that
    \begin{align}\label{E: lower bound on U s=3}
        \E(\Delta_{h\be_2} F_2\cdot {\chi_{2,h}}|\CI(T_2))(x) \geq \delta/2 \quad \textrm{for all}\quad x\in U_h.
    \end{align}
    Hence
    \begin{align}\label{E: original positivity s=3}
        \liminf_{H\to\infty}\E_{h\in[H]} 1_\Lambda(h)\int \Delta_{h\be_2} F_2 \cdot {\chi_{2,h}}\cdot 1_{U_h}\; d\mu > 0.
    \end{align}

    The key complication in this case compared to the one covered by Proposition \ref{P: degree lowering s=2} is the presence of the multiplicative derivative $\Delta_{h\be_2} F_2$ in place of $F_2$. We want to move the differentiation ``inside'' the definition of $F_2$. This step, not present in Proposition \ref{P: degree lowering s=2}, is called \emph{dual-difference interchange}, and we carry it now. 
    
    The rule of thumb in ergodic theory is that if we do not know what to do, it does not hurt to try applying the Cauchy-Schwarz inequality. This is indeed what we will do. First, however, we expand the definition of $T_2^h F_2$ (but not of the second $F_2$), obtaining
    \begin{align*}
        \liminf_{H\to\infty}\E_{h\in[H]}1_\Lambda(h) \lim_{N\to\infty}\E_{n\in[N]}\int F_2\cdot T_2^h\brac{T_2^{-a_2(n)}{f_0}\cdot T_1^{a_1(n)}T_2^{-a_2(n)}{f_1}}\cdot {\chi_{2,h}}\cdot 1_{U_h}\; d\mu > 0.
    \end{align*}
    %Removing the complex conjugates and 
    Moving the (finite) average over $h$ inside the integral and applying the Cauchy-Schwarz inequality, we get
        \begin{align*}
        \liminf_{H\to\infty} \lim_{N\to\infty}\E_{n\in[N]}\norm{\E_{h\in[H]}1_\Lambda(h) \cdot T_2^h\brac{T_2^{-a_2(n)}{f_0}\cdot T_1^{a_1(n)}T_2^{-a_2(n)}{f_1}}\cdot {\chi_{2,h}}\cdot 1_{U_h}}_{L^2(\mu)}^2> 0.
    \end{align*}
    We then expand the square to get the nasty inequality
    \begin{multline*}
         \liminf_{H\to\infty} \limsup_{N\to\infty}\E_{n\in[N]}\E_{h,h'\in[H]}1_{\Lambda^2}(h,h')\\
         \int T_2^h\brac{T_2^{-a_2(n)}{f_0}\cdot T_1^{a_1(n)}T_2^{-a_2(n)}{f_1}}\cdot T_2^{h'}\brac{T_2^{-a_2(n)}\overline{f_0}\cdot T_1^{a_1(n)}T_2^{-a_2(n)}\overline{f_1}}\\ \chi_{2,h}\cdot\overline{\chi_{2,h'}}\cdot 1_{U_h}\cdot 1_{U_{h'}}\; d\mu > 0.
    \end{multline*}

    Because we started with a degree-3 Host-Kra seminorm, the dual-difference interchange required only one application of the Cauchy-Schwarz inequality to perform dual-difference interchange. If we started with the degree-$s$ seminorm, $s-2$ applications would be needed. Given how much notational mess one Cauchy-Schwarz has introduced, we leave it to the reader's vivid imagination to envisage the nightmare that the general case of dual-difference interchange leads to.
    
    Fortunately, we are able to make the notation more palatable by a few manipulations. Composing the integral with $T_2^{-h'}$, defining
    \begin{align*}
        f_{j,h,h'} := \Delta_{(h-h')\be_2}\overline{f_j},\quad \chi_{2,h,h'} := T_2^{-h'}(\chi_{2,h}\cdot\overline{\chi_{2,h'}}),\quad U_{h,h'} := U_h\cap U_{h'},
    \end{align*}
    and swapping the limsup over $N$ with the average over $h,h'$, we get
    \begin{multline*}%\label{E: after dual-difference}
         \liminf_{H\to\infty} \E_{h,h'\in[H]}1_{\Lambda^2}(h,h')\\
         \limsup_{N\to\infty}\abs{\E_{n\in[N]}\int T_2^{-a_2(n)}f_{0,h,h'}\cdot T_1^{a_1(n)}T_2^{-a_2(n)}{f_{1,h,h'}}\cdot \chi_{2,h,h'}\cdot 1_{U_{h,h'}}\; d\mu} > 0.
    \end{multline*}
    
    This expression looks quite like \eqref{E: expanding dual}, and we proceed analogously to the proof of Proposition \ref{P: degree lowering s=2}. Composing the integral with $T_2^{a_2(n)}$ and applying the Cauchy-Schwarz inequality gives us
    \begin{multline*}
        \liminf_{H\to\infty} \E_{h,h'\in[H]}1_{\Lambda^2}(h,h')
         \limsup_{N\to\infty}\norm{\E_{n\in[N]}T_1^{a_1(n)}{f_{1,h,h'}}\cdot T_2^{a_2(n)}(\chi_{2,h,h'}\cdot 1_{U_{h,h'}})}_{L^2(\mu)} > 0.
    \end{multline*}

    The functions $\chi_{2,h,h'}\cdot 1_{U_{h,h'}}$ are nonergodic eigenfunctions of $T_2$, and so invoking Property $(P_1)$ much the same way as in the proof of Proposition \ref{P: degree lowering s=2}, we deduce that 
    \begin{align*}
                \liminf_{H\to\infty} \E_{h,h'\in[H]}1_{\Lambda^2}(h,h')
         \norm{\E(\chi_{2,h,h'}\cdot 1_{U_{h,h'}}|\CI(T_2))}_{L^2(\mu)} > 0.
    \end{align*}
    By the pigeonhole principle, we can find a single\footnote{Technically, this $h'$ may vary with $H$, but we omit this technicality in order not to clutter the notation.} $h'\in\N$ for which
    \begin{align*}
                \liminf_{H\to\infty} \E_{h\in[H]}1_{\Lambda}(h)
         \norm{\E(\chi_{2,h,h'}\cdot 1_{U_{h,h'}}|\CI(T_2))}_{L^2(\mu)} > 0.
    \end{align*}
    Just like in Proposition \ref{P: degree lowering s=2}, we can find $T_2$-invariant sets $U'_h\subseteq U_{h,h'}$ for which
    \begin{align}\label{E: structure on s=3}
        \liminf_{H\to\infty} \E_{h\in[H]}1_{\Lambda}(h)
         \mu(U'_h) > 0,
    \end{align}    
    and such that $\chi_{2,h,h'}$ is $T_2$-invariant on $U'_h$. Given the formula for $\chi_{2,h,h'}$, this means that $\chi_{2,h} = \chi\cdot\lambda_h$ (where we set $\chi:=\chi_{2,h'}$) on $U'_h$ for some $T_2$-invariant $\lambda_h$. Just like in Proposition \ref{P: degree lowering s=2}, we have uncovered some structure on the eigenfunctions $\chi_{2,h}$ although this time the structure is in a sense ``shared'' across different $h$. 

We now combine the lower bound \eqref{E: lower bound on U s=3} and \eqref{E: structure on s=3} with the fact that $U'_h\subseteq U_h$ is $T_2$-invariant to infer that
    \begin{align*}
         \liminf_{H\to\infty} \E_{h\in[H]}1_{\Lambda}(h)\int \E(\Delta_{h\be_2} F_2\cdot {\chi_{2,h}}|\CI(T_2))\cdot 1_{U'_h}\; d\mu > 0.
    \end{align*}
    Using the special structure of  $\chi_{2,h}$ on $U'_h$ and the $T_2$-invariance of $U'_h$, we deduce that
    \begin{align}\label{E: plugging special structure}
                 \liminf_{H\to\infty} \E_{h\in[H]}1_{\Lambda}(h)\int \Delta_{h\be_2} F_2\cdot \chi\cdot \lambda_h\cdot 1_{U'_h}\; d\mu > 0.
    \end{align}
    % By setting $\lambda_h=0$ for $h\notin\Lambda$ (this does not affect any of the discussion above), we get
    % \begin{align*}
    %              \liminf_{H\to\infty} \E_{h\in[H]}\int \Delta_{h\be_2} F_2\cdot \chi\cdot \lambda_h\cdot 1_{U'_h}\; d\mu > 0.
    % \end{align*}
    We invite the reader to compare this expression with \eqref{E: original positivity s=3}, the starting point, to see how arbitrary $T_2$-eigenfunctions $\chi_{2,h}$ have been replaced by the more structured ones. 

    We are almost at the end of the argument; all that remains is to remove the ``low complexity'' terms $1_{\Lambda}(h)\cdot \chi\cdot \lambda_h\cdot 1_{U'_h}$ from the integral. By applying the Cauchy-Schwarz inequality in $h$, we can remove $\chi$, getting
    \begin{align*}
        \liminf_{H\to\infty} \norm{\E_{h\in[H]}T_2^h \overline{F_2}\cdot 1_{\Lambda}(h)\cdot \lambda_h\cdot 1_{U'_h}}_{L^2(\mu)}^2 > 0.
    \end{align*}
    Expanding the square and using the $T_2$-invariance of $1_{\Lambda}(h)\cdot\lambda_h\cdot 1_{U'_h}$, we further infer from the Cauchy-Schwarz inequality that
    \begin{align*}
        \liminf_{H\to\infty} \E_{h,h'\in[H]}\norm{\E(T_2^h F_2\cdot T_2^{h'}\overline{F_2}|\CI(T_2))}_{L^2(\mu)}>0.
    \end{align*}
    A simple exercise shows that this implies $\nnorm{F_2}_{2, T_2}>0,$ as claimed.
    
    Note that if $T_2$ is ergodic, then from \eqref{E: structure on s=3}, we simply get that $\chi_{2,h} = \chi$ for the positive-lower-density set $h\in\Lambda$. Then the conclusion of \eqref{E: plugging special structure} simplifies to
    \begin{align*}
        \liminf_{H\to\infty} \E_{h\in[H]}1_\Lambda(h)\int \Delta_{h\be_2} F_2\cdot \chi\; d\mu %> 0.
        = \liminf_{H\to\infty} \int F_2\cdot \chi \cdot \E_{h\in[H]} T_2^h \overline{F_2}\cdot 1_\Lambda(h)\; d\mu > 0.
    \end{align*}
    The Cauchy-Schwarz inequality gives 
    \begin{align*}
        \lim_{H\to\infty}\norm{\E_{h\in[H]}T_2^h \overline{F_2}\cdot 1_\Lambda(h)}_{L^2(\mu)}>0,
    \end{align*}
    and by expanding the square and pulling the indicators out of the integral, it is easy to conclude that $\nnorm{F_2}_{2, T_2}>0$.
    % The maneuver of defining $\lambda_h = 0$ for $h\notin\Lambda$ has the same effect as taking absolute values outside the integral and using nonnegativity to get
    %     \begin{align*}
    %     \liminf_{H\to\infty} \E_{h\in[H]}\abs{\int \Delta_{h\be_2} F_2\cdot \chi\; d\mu} > 0.
    %     %= \liminf_{H\to\infty} \int F_2\cdot \chi \cdot \E_{h\in[H]} T_2^h \overline{F_2}\; d\mu > 0,
    % \end{align*}
    % and the application of the Cauchy-Schwarz inequality gives
    % \begin{align*}
    %     \lim_{H\to\infty}
    % \end{align*}
\end{proof}

\end{document}
